\section{Introduction}
% In 1987, the concept of `sustainable development' was defined by the United
%  Nations (UN) Brundtland Commission as a development which seeks to ``meet
%  the needs and aspirations of the present without compromising the ability
%  to meet those of the future'' \autocite{WCED1987}.
% Wind turbine manufacturing is an industry shaped by this ideal, with its 
%  end-products contributing towards the UN's 13\textsuperscript{th} sustainable 
%  development goal (SDG) which urges ``urgent action to combat climate change
%  and its impacts'' \autocite{UN2015}.

The early 21\textsuperscript{st} century has seen continued growth in both
 the global human population (from \(6.17\) to \(8.09\)~billion people) and 
 total energy supply required to sustain them (from \(416\) to 
 \(634\)~\unit{\exa\joule}) between 2000 and 2023 
 \autocite{OWID20251027,UN20251027,IEA2025}.
In spite of continued strong reliance upon coal, natural gas, and oil for
 the provision of this demand, concerns regarding global climate and the
 finite nature of fossil reserves have galvanised investment in alternative
 renewable power systems \autocite{IEA2025,IPCC2021}.
Of these renewable, wind power saw the greatest growth in capacity to
 generate electricity (+2.33~\unit{\peta\watt\hour}) ahead of both
 hydroelectric (+1.69~\unit{\peta\watt\hour})
 and solar photovoltaics (+1.60~\unit{\peta\watt\hour}) \autocite{IEA20251027}.
Wind turbine blades are commonly manufactured from a composite 
 comprised of thermosetting polymer matrix and reinforcing material, these
 strong thermoset matrices are chemically cross-linked which makes them
 difficult to reprocess or remould \autocite{Ma2017}.
Given the relatively low rate of recycling and a common rate of 
 replacement (ca. every fifteen years in China \autocite{ChenJ2019}) the
 industry generates a large volume of practically unrecyclable thermoset
 waste which contravenes the UN's 12\textsuperscript{th} SDG to ``ensure
 sustainable consumption and production'' \autocite{UN2015}.
In addressing the industry's sustainability shortcoming identified above, the
 literature is divided into two groups: firstly, papers which focus upon
 improving the recyclability of contemporary materials, whether by thermal,
 chemical, mechanical or high-voltage fragmentation methods
 \autocite{Utekar2020,Beauson2016,ChenR2019}; and secondly, papers aimed at the 
 development of properties in future replacement materials that would 
 facilitate recycling \autocite{Liu2022}, such as introduction of chemically 
 cleavable bonds into the thermoset polymers used \autocite{Tesoro1990}.
The motivation for present work was to ascertain whether Bayesian optimisation
 could be used to sample efficiently improve the degree of polymerisation in
 a novel biodegradable bio-based thermoset.

The composite matrix used for wind turbine blades is comprised of a moulding 
 thermoset, with the most prevalently manufactured pre-thermoset form being 
 general-purpose unsaturated polyester resin (GP-UPR) \autocite{Bobade2016}.
These GP-UPRs are based on a mixture of a reactive diluent (e.g. styrene) and 
 a linear aromatic unsaturated polyester (e.g. poly(propylene glycol maleate
 phthalate)), which are cured chemically or thermally to generate the final
 thermosetting matrix \autocite{Ahamad2001}.
Although polyesters are known to exhibit limited biodegradability (due to 
 their labile ester bonds being prone to both biotic and abiotic hydrolysis
 \autocite{Sharma2011}) it is the case that polyesters with aromatic moieties
 are more resistant than aliphatic counterparts to such degradation
 \autocite{Muller2001}.
Consequently, thermosets derived from GP-UPRs synthesised from the aromatic 
 monomer phthalic acid \autocite{Ahamad2001} can be reasonably expected to
 exhibit a high resistance to biodegradation; which was part of the reason
 for their original choice as a material.
An archetypical example of an aliphatic polyester would be poly(glycerol
 sebacate), which uses the monomers glycerol (triol) and sebacic acid
 (dicarboxylic acid) to form a branching polyester thermoset on thermal
 curing via polyesterification \autocite{Zhang2014}.
Given their natural biodegradability, in the present work we chose to focus 
 upon aliphatic polyesters.

Another concern associated with GP-UPRs with respect to the UN's
 12\textsuperscript{th} SDG is that of a material's provenance.
If the underlying chemistry of a polymer is petrochemical in origin, then
 any carbon liberated from the material at end-of-life (whether by combustion
 for energy or biodegradation in landfill) will release fossil carbon,
 contributing via an anthropogenic route towards global climate change.
Since GP-UPRs are synthesised from the oil-derived
 feedstock monomers propylene glycol, maleic acid, and phthalic acid, the
 resultant thermosets contain endogenic (fossil-derived) carbon
 \autocite{Levi2018}.
In the present work we aim to avoid these feedstocks and focussed instead
 upon a bio-based alternative, with inspiration drawn from a chemical 
 system presented in the literature by Alberts, where a combination of 
 glycerol and citric acid were cured to form a bio-based branched 
 aliphatic polyester thermoset \autocite{Alberts2017}.
We build upon this research by adopting the chemical system and adapting it
 to include a third monomer, itaconic acid, which provides a second
 carboxylic acid with a lower degree of functionality (dicarboxylic acid);
 this second lower-functionality molecule was expected to
 encourage the emergence of alternative thermoset networks.
Notably, Alberts found that physical properties of the end thermoset
 poly(glycerol citrate) could be tuned by modifying the monomeric
 stoichiometries of the prepolymer resin synthesised \autocite{Alberts2017}.
Given this, we identified mass loss as a variable of interest during the
 synthesis of poly(glycerol citrate itaconate) as it is a proxy for
 polymerised content within a sample, and sought to optimise this by 
 tuning the stoichiometries of the prepolymer resins mixed.
Polymerised content could be an indicator of the cross-link density of 
 these samples; a characteristic, which, if maximised, yields improvements
 in mechanical properties such as ultimate tensile strength, Young's modulus,
 and glass transition temperature \autocite{Panic2017}.
Note that the link between sample mass loss and polymerised content is
 predicated on various assumptions: negligible volatilisation of monomers
 during curing, a lack of side-reactions between monomers, and a negligible
 mass of water retained within the final cured matrices of samples.

Poly(glycerol citrate itaconate) is formed from its monomers by
 polyesterification, a process which is dependent upon variables
 such as temperature, stoichiometric ratios, and additive presence 
 \autocite{Alberts2017}.
Additionally, reaction times for rigid solid thermoset materials are on the
 order of weeks when uncatalyzed \autocite{Trenkel2008} (catalysis commonly
 involves strong fuming acids which can corrode equipment \autocite{Shukla2022}).
Given the high dimensionality and temporally expensive function evaluation 
 exhibited by this family of polyesters, a sample efficient global
 optimisation technique, namely Bayesian optimisation, was selected to 
 tackle this problem.
Alternative global optimisation techniques were considered but ruled out;
 for example, the hill-climbing class gradient descent algorithm is reliant on
 derivatives of the objective function which would have
 been uncertain given the noisy function evaluations associated
 with physical systems as opposed to simulated ones \autocite{Stork2022}.
The population class evolutionary and particle swarm optimisation algorithms
 were suitable but avoided in favour of the more sample efficient surrogate
 class of algorithms.
Surrogate class algorithms are variously formulated, but the three major
 forms are distinguished by the surrogate models they leverage: polynomial
 response surface model, kriging (also known as gaussian process regression)
 and radial basis functions \autocite{Roeva2012}.
Kriging is preferable to response surface modelling since it makes no
 assumption of polynomial order, avoiding a priori bias in surrogate models
 generated.
Furthermore, radial basis function-based surrogate models provide no
 estimates of uncertainty around the fitted function; functions and
 associated uncertainties derived using kriging provides the means by which
 Bayesian optimisation effectively balances exploitation and exploration of
 the parameter space.
Hence, for the development of a part-computational and 
 part-laboratory based workflow for the optimisation of mass loss in 
 poly(glycerol citrate itaconate) we decided to use Bayesian optimisation
 predicated on kriging.

Originally coined in the 1970s by Mockus et al. \autocite{Mockus1978}, Bayesian 
 optimisation was later developed further in the 1990s by Jones et al. into a 
 technique called the efficient global optimisation algorithm
 \autocite{Jones1998}- this name and technique was thereafter adopted and
 popularised by the DoE community.
More recently, the original name has seen a revival as the technique and
 several adaptations have been re-discovered and implemented in a range
 of contemporary communities including robotics \autocite{Calandra2016},
 bioengineering \autocite{Bader2023}, and machine learning \autocite{Cho2020}.
In materials science, Bayesian optimisation has been used to predict 
 crystalline structures \autocite{Yamashita2018}, carry out sample-efficient 
 optimisation of density functional theory simulations 
 \autocite{Balachandran2016}, and multi-objective optimisation of stability 
 and ion-conductivity in materials \autocite{Karasuyama2020}.

This paper begins by detailing the polymeric system
 considered, experimental means by which samples were manufactured
 and characterised, as well as the means by which a suitable 
 parameter space was devised.
Hereafter, we detail the algorithmic mechanism behind the Bayesian 
 techniques used before further exploring the initial one-shot quasirandom
 sampling strategy, sequential Bayesian acquisition functions, and 
 one-shot benchmark sampling routines applied during the course of this
 work.
The method concludes with a brief overview of how the physical system
 was computationally emulated such that Bayesian techniques could be
 simulated repeatedly to gain statistical insights.
The results section is split between elucidation of results associated
 with the practical optimisation of the physical system, and those which
 derived from repeated simulation of optimisation routines carried 
 out upon the emulated system.
The paper concludes that the Bayesian sequential routines succeeded in
 improving the mass loss metric from 48\% to 72\% of the theoretical 
 limit, whilst also statistically proving that the `probability of 
 improvement' and `expected improvement' acquisition functions respectively
 exhibited a 1.1\% and 5.3\% probability of discovering values approaching
 the emulated global optimum; this starkly contrasted with the benchmark
 one-shot techniques, for which, it was virtually impossible to discover
 such samples.